\documentclass[12pt,a4paper]{article}
\usepackage[utf8]{inputenc}
\usepackage[T1]{fontenc}
\usepackage[italian]{babel}
\usepackage{amsmath, amssymb, booktabs, tcolorbox, geometry}

\geometry{margin=2.5cm}

\title{Algoritmo DPLL: Analisi dello Scheletro Logico}
\author{Nicola Moscufo}
\date{5 Marzo 2026}

\begin{document}

\maketitle

L'immagine del codice mostra come un computer trasforma la logica proposizionale in un processo sistematico di \textbf{tentativi ed errori}.

\section{La Fase di Obbligo: \texttt{UnitPropagation}}
La prima istruzione operativa è:
\[ (\psi, \mathcal{I}) := \text{UNITPROPAGATION}(\phi, \mathcal{I}') \]

\begin{tcolorbox}[colback=blue!5!white,colframe=blue!75!black,title=Regola 1: La Scelta Forzata]
In questa fase, il computer cerca le \textbf{clausole unitarie}.
\begin{itemize}
    \item Se trova $\{ \lambda \}$, imposta $\lambda = \text{Vero}$.
    \item Se trova $\{ \neg \lambda \}$, imposta $\lambda = \text{Falso}$.
\end{itemize}
\end{tcolorbox}

\section{La Fase di Scelta: \texttt{Splitting Rule}}
Quando non ci sono più clausole unitarie, il programma arriva alla riga:
\begin{center}
    \texttt{else select a literal} $\lambda \in C \in \psi$
\end{center}

\subsection*{Tentativo A: ``Vero''}
\[ \text{if DPLL}(\psi \cup \{\{\lambda\}\}, \mathcal{I}) = \dots \]
\textbf{Significato:} ``Ipotizziamo che $\lambda$ sia Vero.''

\subsection*{Tentativo B: ``Not'' (Backtracking)}
\[ \text{else return DPLL}(\psi \cup \{\{\neg \lambda\}\}, \mathcal{I}) \]
\textbf{Significato:} ``Se il tentativo A fallisce, allora $\lambda$ deve essere Falso.''

\section{Sintesi delle Decisioni}
\begin{table}[h!]
    \centering
    \begin{tabular}{lll}
        \toprule
        \textbf{Situazione} & \textbf{Azione} & \textbf{Logica} \\
        \midrule
        Clausola Unitaria & $\{s\}$ o $\{\neg q\}$ & Decisione Forzata \\
        Bivio (Split) & Scegli $\lambda$ & Decisione Arbitraria \\
        Contraddizione & $\emptyset$ & Backtrack \\
        \bottomrule
    \end{tabular}
\end{table}

\end{document}